\section{An evaluator with weighted dependency paths}\label{sec:evaluation}

We prove the arithmetic evaluation result used below.  Its procedure adds a
scalar multiple of a word's fixed augmented value to one chosen register and
restores every other register, from arbitrary contents on entry.  A caller can
therefore leave its inputs and partial output in a shared pool while a child
borrows the same registers.  The remaining stack cost depends on the
dependency followed: a large loop index is saved only on designated edges.
The translation and interpolation method is related to space-efficient tree
evaluation and its global-storage exposition~\cite{CookMertz2024,Goldreich2024}; all the algebra, register invariants, and
uniformity needed here are proved explicitly.

Throughout this section, $T\ge2$, a word has a common padded width $d\ge1$,
and $r$ is a fixed constant.  A polynomial bound means a bound by $T^C$ for a
fixed constant $C$.  The notation $\polylog T$ permits fixed powers of
$\log(T+2)$, with constants depending on the uniform algorithms under
consideration.  For $v\in\bits^d$, write
\[
 \overline v=(v_1,\ldots,v_d,1).
\]
The last coordinate of an arbitrary vector $z\in\F^{d+1}$ is denoted $z_*$;
it need not equal $1$.

\begin{lemma}[Uniform fields]\label{lem:field}
Given an integer $\Delta\ge1$ bounded polynomially in $T$, one can uniformly
construct a field $\F$ of characteristic two with $|\F|>\Delta$.  Field
elements have $O(\log T)$ bits.  A description of the field, enumeration of its
elements, its arithmetic, and inversion of nonzero elements all require
$\polylog T$ space.
\end{lemma}
\begin{proof}
Start with $\mathbb F_2$.  In a finite field $K$ of characteristic two, the
additive map $x\mapsto x^2+x$ has kernel $\{0,1\}$, since
$x^2+x=x(x+1)$.  Its image therefore contains exactly half of $K$.
Enumerate $a\in K$, and, for each candidate, enumerate $x\in K$ to test
whether $a=x^2+x$.  An $a$ outside the image exists.  The polynomial
$X^2+X+a$ has no root in $K$, hence is irreducible, and
\[
 K[X]/(X^2+X+a)
\]
is a field: the Euclidean algorithm gives an inverse modulo this polynomial
for every nonzero residue of degree less than two.  Its elements are pairs
of elements of $K$ and its cardinality is $|K|^2$.

Repeat until the cardinality first exceeds $\Delta$.  The final cardinality
is at most $\max(2,\Delta^2)$, and its binary vector representation has
$m=O(\log T)$ bits.  The selected coefficients at all levels together use
$O(m)$ bits, since their representation lengths form a geometric sequence.
Arithmetic can be implemented recursively on the pairs, reducing products
using the defining quadratics.  A direct recursive implementation uses at
most a polynomial in $m$ space; its recursion depth is $O(\log m)$.
Enumeration uses an $m$-bit counter.  Alternatively, once multiplication is
available, inversion can be implemented simply by enumerating all elements
until their product with the given nonzero element is $1$.  The searches for
the extension coefficients have the same space bound.  No time bound is
required.
\end{proof}

\begin{lemma}[Homogeneous lookup extensions]\label{lem:homogenization}
For every total function $f:(\bits^d)^a\to\bits$, where $a$ is fixed,
there is a polynomial $\widetilde f$ on $(\F^{d+1})^a$ that agrees with $f$
on augmented Boolean arguments and is homogeneous of degree $d$ separately
in every argument.  If a coordinate of $f$ is computable in
$O(d\polylog T)$ space from its Boolean arguments, its extension can be
evaluated at supplied field vectors in that space, with no recursive word
evaluation.  The same statement applies coordinatewise to augmented
word-valued functions.
\end{lemma}
\begin{proof}
For $u\in\bits^d$ and $z\in\F^{d+1}$, set
\begin{equation}\label{eq:homogeneous-basis}
 B_u(z)=\prod_{\ell=1}^d
 \begin{cases}
  z_\ell,&u_\ell=1,\\
  z_*-z_\ell,&u_\ell=0.
 \end{cases}
\end{equation}
At $z=\overline v$, this is $1$ for $u=v$ and $0$ otherwise.  Thus
\begin{equation}\label{eq:lookup-extension}
 \widetilde f(z^{(1)},\ldots,z^{(a)})
 =\sum_{u^{(1)},\ldots,u^{(a)}\in\bits^d}
 f(u^{(1)},\ldots,u^{(a)})
 \prod_{j=1}^a B_{u^{(j)}}(z^{(j)})
\end{equation}
has the stated agreement and homogeneity.  The zero polynomial is allowed
as a homogeneous polynomial of the indicated degree.  For a word output,
apply the formula to each of its $d$ bit coordinates and to its extra
coordinate, whose Boolean value is $1$.

To evaluate one coordinate, enumerate the $a$ Boolean assignments, compute
the corresponding Boolean function value, and compute the basis products
and running sum using field scalars.  The assignments occupy $ad=O(d)$ bits;
the local Boolean computation occupies $O(d\polylog T)$ bits.  All inputs
are the supplied vectors, read without modification, and immutable
parameters.  This routine makes no
calls to evaluate any word.  Its entire temporary storage can be discarded
when its scalar result has been obtained.
\end{proof}

The augmentation of constant functions is essential.  From
\eqref{eq:homogeneous-basis},
\[
 \sum_{u\in\bits^d}B_u(z)=z_*^d.
\]
Consequently the extension of a globally constant-one one-word function is
$z_*^d$, and that of a globally constant-one two-word function is
$h_*^d a_*^d$.  Neither is replaced by literal $1$.  Moreover, a predicate
that happens to be true for every array value at one particular Boolean
history need not be globally constant.  Such a predicate retains its full
lookup extension.  In particular, the following argument uses these
polynomials even when a shifted augmentation coordinate is zero.

\begin{lemma}[Top-coefficient extraction]\label{lem:top-coefficient}
Let $L$ be an integer with $0\le L<|\F|$, choose any $L+1$ distinct
elements $S_L\subseteq\F$, and set
\[
 \omega_L(s)=\left(\prod_{u\in S_L\setminus\{s\}}(s-u)\right)^{-1}
 \qquad(s\in S_L).
\]
For every univariate polynomial $p$ of degree at most $L$,
\begin{equation}\label{eq:top-coefficient}
 [X^L]p(X)=\sum_{s\in S_L}\omega_L(s)p(s).
\end{equation}
If $Q$ is homogeneous of degree $L$, then, for arbitrary field vectors $a,v$,
\begin{equation}\label{eq:translation-extraction}
 \sum_{s\in S_L}\omega_L(s)Q(a+sv)=Q(v).
\end{equation}
Choosing $S_L$ as the first $L+1$ elements in the uniform field enumeration,
its points and weights can be enumerated and recomputed in $\polylog T$
space when $L$ is polynomially bounded in $T$.
\end{lemma}
\begin{proof}
The polynomials
\[
 \prod_{u\in S_L\setminus\{s\}}\frac{X-u}{s-u}
\]
interpolate the values at $S_L$.  Subtracting the resulting interpolant from
$p$ gives a polynomial of degree at most $L$ with $L+1$ distinct roots, so
the difference is zero.  Taking the coefficient of $X^L$ proves
\eqref{eq:top-coefficient}.  For each monomial of total degree $L$, the
coefficient of $X^L$ after substituting $a+Xv$ is the same monomial in $v$.
Summing proves \eqref{eq:translation-extraction}.  Enumeration of the first
$L+1$ field elements and a running product compute each weight using only
field scalars and counters.  All denominators are nonzero.  This proof uses
neither factorials nor derivatives, and applies in characteristic two even
when $L$ exceeds the characteristic.
\end{proof}

We now specify the implicit computation to which the evaluator applies.
There is a finite acyclic directed graph of word identifiers, with edges
from a word to the words on which it depends.  The graph has polynomially
many vertices, and identifiers use $O(\log T)$ bits.  Every vertex $V$ has a
fixed semantic value $v(V)\in\bits^d$, determined by the input and immutable
parameters, independently of all working registers.  There are three types
of vertices:
\begin{enumerate}
\item A base word has a locally computable value.
\item An ordinary vertex has a rule
\[
 v(V)=f_V\bigl(v(U_1),\ldots,v(U_a)\bigr),\qquad 1\le a\le r,
\]
where $f_V$ is a total Boolean word function.
\item A history vertex has a distinguished dependency $H$, further
 dependencies $A_1,\ldots,A_N$, and an integer $0\le g\le d$.
 There are total Boolean functions
\[
 G_Y:\bits^d\longrightarrow\bits^d,
 \qquad
 P_{Y,i}:\bits^d\times\bits^d\longrightarrow\bits
 \quad(Y\in\bits^g, 1\le i\le N),
\]
where the vertex identifier is an additional suppressed parameter.
At the semantic arguments they satisfy the coordinatewise identity
\begin{equation}\label{eq:abstract-history}
 \overline{v(V)}=
 \sum_{Y\in\bits^g}
 \overline{G_Y(v(H))}\prod_{i=1}^N
 P_{Y,i}\bigl(v(H),v(A_i)\bigr)
\end{equation}
in characteristic two.  It suffices, in particular, that exactly one $Y$
has all predicates equal to one and that its output is $v(V)$.
\end{enumerate}
The identity \eqref{eq:abstract-history} is required only at the semantic
arguments; arbitrary Boolean tuples need not describe consistent histories.

Here and below $d,N$ have uniform polynomial upper bounds.  Given a vertex
identifier, its type, local parameters, and any indexed dependency can be
computed using $\polylog T$ space.  Base coordinates and coordinates of the
Boolean functions $f_V,G_Y,P_{Y,i}$ are uniformly computable in
$O(d\polylog T)$ space from their explicit Boolean arguments, the vertex
identifier, and, where applicable, $Y,i$.  These local routines make no word
evaluation calls.  A single supplied polynomial upper bound $\Delta$ covers
$d$ and every $(N+1)d$.  Additional immutable data can be read by these
algorithms; if such data occupy work space, that storage is counted
separately, once.  These assumptions make navigation and local computation
part of the lemma, rather than treating an exponentially large truth table
as a freely available input.

Give ordinary dependency edges weight $1$.  At a history vertex, give the
edge to $H$ weight $1$, and each edge to $A_i$ weight $g+1$.  Dependencies
with the same identifier but different roles are treated as separate edges.
Let
\begin{equation}\label{eq:weighted-path-definition}
 W(V)=\max_{\pi}\sum_{e\in\pi}\wt(e),
\end{equation}
where $\pi$ ranges over dependency paths starting at $V$ and ending at a
base word.  A base word has $W(V)=0$.

\begin{lemma}[Evaluation along weighted paths]\label{lem:weighted-evaluation}
Under the preceding assumptions, the augmented value of a word $V$ can be
computed deterministically in
\begin{equation}\label{eq:weighted-evaluation-space}
 O\bigl((d+W(V)+1)\polylog T\bigr)
\end{equation}
work space, in addition to separately stored immutable auxiliary data.
More strongly, a pool of exactly $p=\max(3,r+1)$ registers in
$\F^{d+1}$ supports a procedure $\Add(V,\lambda,t)$, for every
$\lambda\in\F$ and target index $t$, with the following contract for
\emph{every} initial assignment to the entire pool:
\begin{equation}\label{eq:add-contract}
 R_t^{\mathrm{out}}=R_t^{\mathrm{in}}+\lambda\overline{v(V)},
 \qquad
 R_j^{\mathrm{out}}=R_j^{\mathrm{in}}\quad(j\ne t).
\end{equation}
The space bound \eqref{eq:weighted-evaluation-space} includes this shared
pool and all temporary and suspended storage.
\end{lemma}
\begin{proof}
Use the field from Lemma~\ref{lem:field} for the supplied degree bound
$\Delta$.  Each register occupies $O(d\log T)$ bits.  We prove the contract
by induction in the acyclic dependency graph, simultaneously for all
scalars, targets, and initial register assignments.  At every recursive
call, the scalar request and the word identifier are ordinary saved data;
neither is encoded by a register whose contents the call can change.

\paragraph{Base words.}
Generate one semantic bit at a time and add its product with $\lambda$ to
the corresponding target coordinate.  Add $\lambda$ to the extra
coordinate.  No other register changes.  Local generation uses the stated
scratch space and no recursive word calls.

\paragraph{Ordinary vertices.}
Let $U_1,\ldots,U_a$ be the dependencies.  Choose distinct slots
$q_1,\ldots,q_a$, all different from $t$; there are enough slots since
$p\ge r+1$.  Apply Lemma~\ref{lem:homogenization} to each coordinate of
$\overline{f_V}$, including its extra coordinate.  Denote the resulting
vector of polynomials by $\widetilde f_V$.

Nest $a$ interpolation loops.  At level $j$, for each $s_j\in S_d$, execute
\[
 \Add(U_j,s_j,q_j),\qquad
 \text{the next loop level},\qquad
 \Add(U_j,-s_j,q_j).
\]
At the innermost level, compute the coordinates of
$\widetilde f_V(R_{q_1},\ldots,R_{q_a})$ one at a time, and add each scalar
to the corresponding coordinate of $R_t$ with coefficient
$\lambda\prod_{j=1}^a\omega_d(s_j)$.  Discard each local lookup
computation's scratch after obtaining its scalar result, before any further
word call.

The induction hypothesis implies that each forward shift translates the
chosen slot by $s_j\overline{v(U_j)}$, and that each reverse shift restores
its baseline.  Every other slot has its current value restored after each
call.  In particular, at a lookup computation the arguments are
$a_j+s_j\overline{v(U_j)}$, where $a_j$ are their values on entry to this
invocation.  Applying \eqref{eq:translation-extraction} successively to the
separately homogeneous arguments shows that the total increment of the
target is
\[
 \lambda\widetilde f_V
    (\overline{v(U_1)},\ldots,\overline{v(U_a)})
 =\lambda\overline{v(V)}.
\]
Every non-target slot is restored.  The nesting depth inside this invocation
is at most the constant $r$, and only field scalars, counters, slot indices,
and identifiers are saved at its recursive word calls.

\paragraph{History vertices: the polynomial.}
Write $h_0=\overline{v(H)}$ and $a_i=\overline{v(A_i)}$.  Extend each
coordinate of $\overline{G_Y}$ by the one-word lookup formula and each
$P_{Y,i}$ by the two-word formula, calling the results
$\widetilde G_Y$ and $\widetilde P_{Y,i}$.  Set
\begin{equation}\label{eq:history-polynomial}
 F(h;a_1,\ldots,a_N)
 =\sum_{Y\in\bits^g}\widetilde G_Y(h)
       \prod_{i=1}^N\widetilde P_{Y,i}(h,a_i).
\end{equation}
It is homogeneous of degree $D=(N+1)d$ in $h$, and each factor is
homogeneous of degree $d$ in its own array argument.  In particular the
extra coordinate of $\widetilde G_Y(h)$ is $h_*^d$.
Equation~\eqref{eq:abstract-history} and Boolean agreement imply
\[
 F(h_0;a_1,\ldots,a_N)=\overline{v(V)}.
\]
All interpolation degrees $d$ and $D$ are at most $\Delta<|\F|$.

\paragraph{History vertices: the schedule.}
Choose slots $q_H,q_A,t$ distinct.  Let $h$ and $a$ denote the entry contents
of $q_H,q_A$.  Indentation in the following schedule specifies the loop
bodies.  Each lookup finishes before its scratch is released.

\begin{small}
\begin{tabbing}
\hspace{1.4em}\=\hspace{1.4em}\=\hspace{1.4em}\=\hspace{1.4em}\=\hspace{1.4em}\=\kill
\textbf{for} $s\in S_D$:\\
\> $\Add(H,s,q_H)$ \quad (no coordinate or $Y$ loop is active)\\
\> \textbf{for} $z\in\{1,\ldots,d,*\}$:\\
\>\> \textbf{for} $Y\in\bits^g$:\\
\>\>\> $u\leftarrow\widetilde G_{Y,z}(R_{q_H})$ \quad (local lookup)\\
\>\>\> \textit{release lookup scratch, keeping the scalar $u$}\\
\>\>\> \textbf{for} $i=1,\ldots,N$:\\
\>\>\>\> $b_i\leftarrow0$\\
\>\>\>\> \textbf{for} $s'\in S_d$:\\
\>\>\>\>\> $\Add(A_i,s',q_A)$\\
\>\>\>\>\> $c\leftarrow\widetilde P_{Y,i}(R_{q_H},R_{q_A})$ \quad (local lookup)\\
\>\>\>\>\> \textit{release lookup scratch, keeping scalar $c$ in the frame}\\
\>\>\>\>\> $\Add(A_i,-s',q_A)$ \quad (retain $c$ across this call)\\
\>\>\>\>\> $b_i\leftarrow b_i+\omega_d(s')c$\\
\>\>\>\> $u\leftarrow u b_i$; \textit{ discard $b_i$}\\
\>\>\> $R_{t,z}\leftarrow R_{t,z}+\lambda\omega_D(s)u$\\
\> \textit{discard data from the coordinate, $Y$, and factor loops}\\
\> $\Add(H,-s,q_H)$ \quad (no coordinate or $Y$ loop is active)
\end{tabbing}
\end{small}

Each completed factor is multiplied into $u$ and discarded, so no list of
factors is stored.  Both calls to $H$ occur outside the coordinate and $Y$
loops.  All loops may take arbitrarily long, but their indices have the
stated finite lengths.

For fixed $s$, the induction hypothesis restores $R_{q_H}=h+sh_0$ after
each array call and restores $R_{q_A}=a$ after each inner sample.  Hence
\eqref{eq:translation-extraction}, applied only to the array argument,
gives
\[
 b_i=\sum_{s'\in S_d}\omega_d(s')
       \widetilde P_{Y,i}(h+sh_0,a+s'a_i)
     =\widetilde P_{Y,i}(h+sh_0,a_i).
\]
The coordinate and $Y$ loops therefore add
$\lambda\omega_D(s)F(h+sh_0;a_1,\ldots,a_N)$ to the target.  Outer
coefficient extraction adds exactly
$\lambda F(h_0;a_1,\ldots,a_N)=\lambda\overline{v(V)}$ in total.
The final reverse shift restores $q_H$, and the other non-target slots are
also restored.  This proves the contract for a history vertex.  No
independence assumption among the semantic array values $a_i$ is involved:
each coefficient identity holds for their fixed values.  Their independence
from the mutable register pool is the required condition.

\paragraph{Why one pool suffices.}
Distinctness of target and input slots is imposed within the current
invocation.  A child invocation may use any physical slots, including the
caller's partially accumulated target, the caller's shifted history, or an
ancestor's input slots.  The induction hypothesis quantifies over every
entry assignment, so the child restores all such slots other than its own
target before returning.  The caller uses those values only after the child
returns.  A reverse shift is another invocation of the same contract and
restores its target by subtraction of a fixed semantic vector; it does not
require other registers to have the values they had before the forward
shift.  Thus no ancestor needs to save a vector, and no private vector pool
is allocated on descent.  In characteristic two the negative scalar equals
the positive scalar, which still performs the required cancellation.

\paragraph{Scratch lifetimes and the stack.}
The vector pool is global.  So is one reusable area of
$O(d\polylog T)$ bits for local Boolean assignments, replay work, and lookup
summations.  Every use of this area ends before a recursive word call.
Only its scalar result, when needed, is moved to the saved frame; no local
Boolean assignment survives such a call.  This temporary area need not be
restored, since it has no live caller data when a child starts using it.
The active $Y$ index has $g$ bits and is distinct from this disposable lookup
scratch.

For the saved data use a contiguous stack with a top pointer measuring its
occupied length.  Arrange an invocation's short outer loop data first and its
currently live inner loop data above them.  Immediately before a word call,
set the top just beyond the caller data that must survive that call; the child
allocates its frames from that position and returns the top to it.  In
particular, after the coordinate
and $Y$ loops end, rewind the top to the end of the short outer frame before
the reverse call to $H$.  That call may overwrite the released loop cells.
Thus sequential uses of those cells contribute their maximum extent, not
their sum, to the work-space bound.  The pointer itself has $O(\log T)$ bits
because the graph and word widths are polynomially bounded.

On an ordinary edge, the suspended frame contains a constant number of
point indices and weights, field scalars, identifiers, slot indices, and
control-state flags: $\polylog T$ bits in all.  The same bound holds on a
history-to-$H$ edge, because both calls to $H$ occur outside the coordinate
and $Y$ loops.  On a history-to-$A_i$ edge, one additionally retains the
current $Y$, along with its coordinate and factor indices, the product $u$,
one factor accumulator, and the inner interpolation state.  This occupies
$O((g+1)\polylog T)$ bits.  Forward and reverse calls have the same bound
and occur sequentially.

At any instant, active word calls form a dependency path.  Summing the
bounds on their suspended frames gives $O(W(V)\polylog T)$ bits.  The
current invocation has at most $g\le d$ bits of active $Y$ data in addition
to its short bookkeeping.  Adding the pool, the maximum local scratch,
field description, and current frame proves
\eqref{eq:weighted-evaluation-space}.  Navigation is uniform by hypothesis,
and the graph is acyclic, so all recursive calls terminate.  Finally,
initialize a target register to zero and call $\Add(V,1,t)$ to obtain the
augmented word itself.
\end{proof}
